\documentclass[10pt,a4paper]{article}

\usepackage[left=15mm,right=15mm,top=13mm,bottom=16mm]{geometry}
\usepackage{fontspec}

\setmainfont{TeX Gyre Heros}

\usepackage{academicapp}

\hypersetup{pdftitle={Shiqi Zhang Academic Curriculum Vitae}}

\setstretch{1.00}

\setlength{\parskip}{1.6pt}

\SetAppFooter{Academic Curriculum Vitae | Updated November 2026}

\begin{document}

\ApplicantName{\textbf{SHIQI ZHANG, Ph.D.}}\vspace{1pt}

\ApplicantSubtitle{\textbf{Academic Curriculum Vitae | Power Electronics • Converter-Dominated Systems • Renewable Hydrogen}}\vspace{1pt}

\ApplicantMeta{Stockholm, Sweden  |  shizha@kth.se  |  +86 188 3388 4449  |  https://s-q-z.github.io/\\ Google Scholar: f6Bvb1UAAAAJ}\vspace{2pt}

\section*{Research profile}

Power-electronics researcher developing reliable and intelligent conversion systems for renewable-energy, storage, and hydrogen applications. Research integrates converter topology and modulation, impedance-based stability, grid-forming and weak-grid control, fault diagnosis and fault-tolerant operation, and real-time/prototype validation.\par

\section*{Academic appointment}

\textbf{Researcher} | KTH Royal Institute of Technology, Department of Electric Power and Energy Systems | Stockholm, Sweden | 2026–present\par

\begin{itemize}\item Power-electronics-enabled renewable and hydrogen systems; converter-grid/electrolyzer interactions; control-oriented modeling; hardware-oriented validation.\end{itemize}

\begin{itemize}\item Simulation and validation workflows integrating physics-based models, advanced control, intelligent optimization, and real-time experimentation.\end{itemize}

\section*{Education}

\textbf{Ph.D. in Electrical Engineering} | Yanshan University, China | January 2026\par

\begin{itemize}\item GPA 90.4/100; rank 1/54. Research in converter control and modulation, renewable-integrated converters, weak-grid stability, impedance modeling, fault diagnosis, and fault-tolerant operation.\end{itemize}

\textbf{M.Eng. in Emerging Power Systems} | Curtin University, Australia | August 2022\par

\begin{itemize}\item GPA 4.26/5.0; rank 1/35. Coursework included Distributed Generation, Smart Grid Technologies, and Power System Protection and Control.\end{itemize}

\section*{Research leadership and funded-project experience}

\textbf{Innovation Project of Hebei Province} | Principal Investigator | October 2023–March 2025\par

\begin{itemize}\item Led modeling and control research for converter-interfaced renewable-energy systems under intermittent, weak-grid, and disturbance-prone conditions.\end{itemize}

\begin{itemize}\item Developed multi-timescale models and constraint-aware learning-based energy management; validated through simulation, real-time simulation, and experiments.\end{itemize}

\textbf{National Key R\&D Program – Major Energy Project} | Participant | January 2021–January 2026\par

\begin{itemize}\item Contributed to coordinated control, ripple suppression, reliability enhancement, and fault-tolerant operation of renewable-hydrogen converter systems.\end{itemize}

\section*{Research contributions}

\begin{itemize}\item Low-ripple and high-step-down conversion: converter and modulation methods for electrolyzers and multi-port renewable-hydrogen systems.\end{itemize}

\begin{itemize}\item Weak-grid and grid-interactive stability: PLL design, impedance modeling, and converter-grid interaction analysis.\end{itemize}

\begin{itemize}\item Fault-aware conversion: interpretable data-driven diagnosis, reconfiguration, and capability-preserving control.\end{itemize}

\begin{itemize}\item System integration: coordinated PV, battery, electrolyzer, DC-link, and supervisory operation across timescales.\end{itemize}

\section*{Selected technical projects}

\textbf{Three-level energy-storage inverter topology and control} | Industry collaboration with a global inverter manufacturer | May--November 2023. Evaluated three-level converter structures for grid-connected storage, developed the associated control design, and compared operating performance to guide topology and control optimization for industrial inverter applications.\par

\textbf{Zhangbei wind-power grid-integration project} | Industry collaboration with Longyuan Renewable Energy | March--October 2023. Conducted converter--grid interaction and impedance-based stability analysis for wind-farm integration, then developed oscillation-suppression and impedance-reshaping control measures for weak-grid operation. The work contributed to granted invention patent CN116544963B, ``Impedance reshaping control method for wind turbines'' (co-inventor; granted 24 October 2023).\par

Constrained reinforcement learning for multi-port modulation. Developed a DDPG framework for a triple-active-bridge system with five phase-shift decision variables, feasibility-preserving action mapping, target networks, replay, and randomized 200–300 A operating references. The project connects learning-based optimization to explicit converter constraints and low-ripple objectives.\par

Fault diagnosis and fault-tolerant operation. Combined SSA-optimized random forests for open-circuit fault identification with auxiliary-switch reconfiguration and power sharing. Disseminated through IEEE APEC 2025 and IEEE Transactions on Power Electronics (2025).\par

Grid-forming electrolyzer interface. Built switching-level models of a six-IGBT rectifier with PWM and DC-link dynamics; examined partial-load operation, ramps, undervoltage behavior, delayed restart, and staged recovery.\par

Multi-stack renewable-hydrogen scheduling. Studied coordinated operation of fifteen 1 MW PEM electrolyzer stacks using a 30-day NREL photovoltaic sequence and operational constraints.\par

\section*{Experimental and computational capabilities}

MATLAB/Simulink • PLECS • PSCAD • Python • RT-LAB/OP5707 • DSP • FPGA • HIL • switching and averaged models • impedance analysis • prototype testing • optimization • machine learning\par

\section*{Teaching, mentoring, and academic service}

\begin{itemize}\item Mentored graduate researchers in converter modeling, MATLAB/Simulink, experimental planning, data analysis, manuscript development, and validation.\end{itemize}

\begin{itemize}\item Supported laboratory-oriented training in converter testing, HIL, DSP/FPGA control, and prototype commissioning.\end{itemize}

\begin{itemize}\item Peer reviewer for Nature Communications, IEEE Transactions on Power Electronics, IEEE Transactions on Transportation Electrification, IEEE JESTPE, IET Power Electronics, and Journal of Energy Storage.\end{itemize}

\begin{itemize}\item Oral presentations at IEEE APEC 2025, IEEE PEDS 2023, IEEE PEAC 2022, and SPIES 2022.\end{itemize}

\section*{Awards and honors}

\begin{itemize}\item First Prize, Science and Technology Progress Award, China Power Supply Society, 2025 (fourth contributor).\end{itemize}

\begin{itemize}\item China Telecom Tianyi Scholarship, 2025 (one of 98 nationwide awardees).\end{itemize}

\begin{itemize}\item National Scholarship for Graduate Students, 2023, 2024, and 2025.\end{itemize}

\begin{itemize}\item Youth Medal of Hebei Province, 2025.\end{itemize}

\begin{itemize}\item National Gold Awards, China International College Students Innovation and Entrepreneurship Competition, 2021 and 2022.\end{itemize}

\begin{itemize}\item Silver Medal, 49th Geneva International Exhibition of Inventions, 2024 (third contributor).\end{itemize}

\section*{Peer-reviewed journal articles}

[J1] S. Zhang et al., “Fault Diagnosis and Fault-Tolerant Operation Strategy for Triple-Port Hydrogen Production System,” IEEE Transactions on Power Electronics, 40(8), 11526–11536, 2025.\par

[J2] S. Zhang et al., “An Industrial Overview of Variable-Frequency Drives for Hydrogen Compressors in Hydrogen Transportation,” Renewable and Sustainable Energy Reviews, 223, 116009, 2025.\par

[J3] S. Zhang et al., “Low-Ripple Modulation Strategy for a Photovoltaic-Based Triple-Port Hydrogen Production System,” IEEE Transactions on Industrial Electronics, 73(5), 7084–7095, 2026.\par

[J4] S. Zhang et al., “Overview of US Patents for Energy Management of Renewable Energy Systems with Hydrogen,” International Journal of Hydrogen Energy, 48(26), 9574–9591, 2023.\par

[J5] X. Guo, S. Zhang (corresponding author) et al., “Advancements in Photovoltaic Electrolysis for Green Hydrogen Production: A Comprehensive Review and Comparative Analysis of Modeling Approaches,” IEEE Transactions on Power Electronics, 40(7), 10000–10026, 2025.\par

[J6] X. Guo, S. Zhang (corresponding author) et al., “A Novel DC–DC Converter for Electrolyzer With Low Ripple and High Step Down,” IEEE Transactions on Industrial Electronics, 71(10), 12476–12486, 2024.\par

[J7] X. Guo, S. Zhang et al., “A New Multi-Mode Fault-Tolerant Operation Control Strategy of Multiphase Stacked Interleaved Buck Converter for Green Hydrogen Production,” International Journal of Hydrogen Energy, 47(71), 30359–30370, 2022.\par

[J8] X. Guo, S. Zhang (corresponding author) et al., “A Novel Phase-Locked Loop Structure to Enhance Converter Stability in Weak Grids,” IEEE Transactions on Power Electronics, 38(11), 13855–13865, 2023.\par

[J9] X. Guo, F. Ding, S. Zhang (corresponding author) et al., “Enhanced Green Hydrogen Production Using a Novel Converter with Extremely Low Current Ripple and High Step-Down Ratio,” Fuel, 398, 135424, 2025.\par

\section*{Peer-reviewed journal articles (continued)}

[J10] X. Guo, H. Zhu, and S. Zhang (corresponding author), “Overview of Electrolyser and Hydrogen Production Power Supply from Industrial Perspective,” International Journal of Hydrogen Energy, 49, 1048–1059, 2024.\par

[J11] W. Fang, Y. Teng, S. Zhang (corresponding author) et al., “Impedance Modeling and Stability Analysis of Electrolysis System for Hydrogen Production under Weak Grid,” Fuel, 374, 132403, 2024.\par

[J12] X. Guo, X. Hu, and S. Zhang (corresponding author), “Application Status of Variable-Frequency Drive in Hydrogen Fuel Cell Air Compressors from an Industrial Viewpoint: A Review,” Sustainable Energy Technologies and Assessments, 64, 103716, 2024.\par

[J13] A. Syed et al., “Next-Generation Control for Electrolyzers: A Review of GPT-Based AI Frameworks in Renewable Hydrogen Systems,” Applied Energy, 406, 127300, 2026.\par

[J14] A. Syed et al., “A Novel High Step-Down DC–DC Converter With Low Ripple for Electrolyzer Applications,” IEEE Transactions on Industrial Electronics, 72(12), 13405–13417, 2025.\par

[J15] X. Zhang et al., “An Overlap-Time Optimization SVM Method Based on Vector Sequential Reconstruction for Current Source Converters,” IEEE Journal of Emerging and Selected Topics in Power Electronics, Early Access, 2025.\par

[J16] Y. Teng et al., “Hydrogen in Practice: A Global Review of Demonstration Projects from Production to Integration,” Renewable and Sustainable Energy Reviews, 242, 117314, 2026.\par

[J17] Y. Gao, S. Zhang et al., “Optimal Dispatch of Hydrogen Production Systems: Integrating Electrolyzer Start-Up Optimization and Deep Reinforcement Learning,” Energy, 360, 141612, 2026.\par

\section*{Peer-reviewed conference papers}

[C1] S. Zhang et al., “Fault Diagnosis and Tolerant Strategy for Triple-Port Hydrogen Converter Using SSA-Optimized Random Forest Algorithm,” IEEE APEC, 2025, pp. 757–760. Oral presentation.\par

[C2] X. Guo, S. Zhang (corresponding author) et al., “A Novel DC–DC Converter for Electrolyzer with Low Ripple and High Step-Down Ratio,” IEEE PEDS, 2023, pp. 1–5. Oral presentation.\par

[C3] S. Zhang and X. Guo, “Optimal Configuration of Integrated Target Distributed Power Generation Based on Improved JAYA Algorithm and AHP,” IEEE PEAC, 2022, pp. 684–687. Oral presentation.\par

[C4] S. Zhang et al., “A Novel Three-Phase Multilevel Converter for Medium-Voltage Application,” SPIES, 2022, pp. 294–297. Oral presentation.\par

\section*{Patents and intellectual property}

35 patent records, including 11 as first inventor. Selected granted and pending records:\par

\begin{itemize}\item Hydrogen production converter and control method, CN118174544B, granted 2024 (first inventor).\end{itemize}

\begin{itemize}\item Energy management method for microgrids with hydrogen energy storage, CN115085229B, granted 2022 (first inventor).\end{itemize}

\begin{itemize}\item Fault diagnosis and recovery control method for interleaved Buck circuits, CN115065243B, granted 2022 (first inventor).\end{itemize}

\begin{itemize}\item Control and fault-tolerant method for multiphase interleaved parallel converters, CN115347774B, granted 2023 (first inventor).\end{itemize}

\begin{itemize}\item Improved transformerless three-phase multilevel inverter and control method, CN115987127B, granted 2023 (first inventor).\end{itemize}

\begin{itemize}\item Hydrogen production power converter, CN117040267B, granted 2024 (first inventor).\end{itemize}

\begin{itemize}\item Method for constructing a model of an electrolysis hydrogen production system, CN117037928B, granted 2025 (first inventor).\end{itemize}

\begin{itemize}\item Intelligent optimization control method for satellite power fault reconstruction, CN115079571B, granted 2022 (first inventor).\end{itemize}

\begin{itemize}\item Multi-mode dual-layer fault diagnosis method for fuel cell systems, CN113962259B, granted 2022 (first inventor).\end{itemize}

\begin{itemize}\item Energy management method for renewable energy systems, CN115411770B, granted 2023 (first inventor).\end{itemize}

\begin{itemize}\item Multi-objective distributed generation optimal configuration method, CN113890094B, granted 2022 (first inventor).\end{itemize}

\begin{itemize}\item Impedance modeling method for three-phase four-leg grid-forming converters, CN118364766B, granted 2025 (co-inventor).\end{itemize}

\begin{itemize}\item Energy management and optimal scheduling method for renewable energy systems, CN118449171B, granted 2025 (co-inventor).\end{itemize}

\begin{itemize}\item Impedance reshaping control method for wind turbines, CN116544963B, granted 24 October 2023 (co-inventor; linked to the Zhangbei wind-power grid-integration project).\end{itemize}

\end{document}
